% !PRINTED 580
\weberpage{580}
%<<CONT
und lassen sich in folgender Weise entwickeln:
\weq{17}{\frac{\sum\limits_{-\infty,\infty}^{\nu} q^{\left(\nu+\frac12+\frac{h'}{m}\right)^2} e^{\frac{(2\nu+1)\pi i}{m}h}}{\sum\limits_{-\infty,\infty}^{\nu} q^{\left(\nu+\frac{h'}{m}\right)^2} e^{\frac{2\nu\pi i}{m}h}}.}
Im Nenner dieses Ausdruckes kommt nur ein Glied mit
niedrigster Potenz von $q$ vor, das man erhält, wenn man für $\nu$
die dem Bruch $-h'/m$ zunächst gelegene ganze Zahl setzt.

Im Zähler kommt im allgemeinen auch nur ein niedrigstes
Glied vor, dessen Koeffizient eine Einheit ist. In dem besonderen
Falle, wo $h'/m$ eine ganze Zahl ist, kommen aber im Zähler zwei
gleiche niedrigste Glieder vor, nämlich (für $h'=0$) $\nu=0,-1$.
Diese geben zusammen
\[
q^{\frac14}\,2\cos\frac{h\pi}{m}
\]
und $2\cos\dfrac{h\pi}{m}$ ist eine algebraische Einheit. [Zu schließen aus
Bd. I, § 144, (19).]

Demnach läßt sich sowohl der Ausdruck (17), \gesperrt{als auch sein
reziproker Wert} nach steigenden Potenzen von $q$ in der Art
entwickeln, daß der Koeffizient des ersten Gliedes den Wert 1
hat und alle übrigen Koeffizienten ganze algebraische Zahlen sind.
% [source sic: the preceding exceptional case gives the generally nontrivial algebraic unit 2 cos(h\pi/m)]
Daraus schließt man, ebenso wie in bezug auf $A(x)$, daß sich
$B(x)\,C(x)$ in der Weise darstellen läßt:
\weq{18}{B(x)\,C(x) = x^{2m^2-2} + b_1 x^{2m^2-4} + \cdots + 1,}
worin die $b_1, b_2 \ldots$ ganze rationale Funktionen von $\lambda$ mit ganzen
rationalen Zahlenkoeffizienten sind. Außerdem sind die Koeffizienten
 $b$, die gleichweit vom Anfang und vom Ende abstehen,
einander gleich.

Betrachten wir noch die aus (8) fließende Gleichung:
\weq{19}{x A(x) - \sqrt{\varkappa}\,\mathrm{sn}\,mv\,D(x) = 0,}
die in bezug auf $x$ vom Grade $m^2$ ist. Ihre Wurzeln sind die
Größen:
\[
\sqrt{\varkappa}\,\mathrm{sn}\,v, \quad \sqrt{\varkappa}\,\mathrm{sn}(v+\Omega_{h,h'}).
\]
Das Produkt dieser Wurzeln ist, vom Vorzeichen abgesehen,
gleich dem unabhängigen Gliede in dieser Gleichung, also gleich
\[
\pm \sqrt{\varkappa}\,\mathrm{sn}\,mv,
\]
und daraus ergibt sich:
\weq{20}{\pm\frac{\sqrt{\varkappa}\,\mathrm{sn}\,mv}{\sqrt{\varkappa}\,\mathrm{sn}\,v} = \prod^{h,h'} \sqrt{\varkappa}\,\mathrm{sn}(v+\Omega_{h,h'}).}
